% ============================================================
%  GENERAL CONCLUSION AND PERSPECTIVES
%  Save as: Chapters/conclusion.tex
%  Input in main.tex after chapter4
% ============================================================

\chapter*{General Conclusion and Perspectives}
\addcontentsline{toc}{chapter}{General Conclusion and Perspectives}
\label{chap:conclusion}

\section*{Synthesis of Contributions}

This thesis addressed the problem of optimising federated
learning for privacy-constrained healthcare applications,
validated on breast DCE-MRI molecular subtype classification
--- a domain where legal constraints under Law~18-07 of June
10, 2018 and international data protection frameworks
prohibit direct patient data sharing across institutions,
and where the rarest subtype (HER2-enriched, 6.2\% of cases)
represents only 46 patients in the available dataset.
Two complementary research threads were pursued: a
centralised baseline establishing the performance ceiling,
and a federated optimisation study quantifying how closely
a privacy-preserving distributed system can approach it.

\textbf{Centralised contributions.}
The ConvNeXt-Nano backbone with Gated Attention Multiple
Instance Learning, trained through a three-stage pipeline
combining Label-Distribution-Aware Margin loss, Adaptive
Sharpness-Aware Minimization, Stochastic Weight Averaging,
and Classifier Retraining, achieves four-class ensemble
macro F1 = 0.342 and binary macro F1 = 0.667 without
tumour segmentation masks.
These results are competitive with published mask-guided
methods when evaluated under equivalent constraints:
whole-breast analysis, single publicly available cohort
(737 patients), and patient-level cross-validation that
eliminates the optimistic bias of sample-level splits.
The three-stage training pipeline addresses the 10.4:1
class imbalance systematically: LDAM corrects the training
signal, ASAM finds flatter loss basins that generalise
across fold distributions, SWA stabilises the trajectory,
and cRT recalibrates the head after backbone fine-tuning.
The Stage~2 to cRT gain of 0.064 F1 points in the binary
setting confirms that decoupling representation learning
from head calibration is effective even in a centralised
context.

The radiomics fusion experiment yielded a negative result
(F1 decreased from 0.602 to 0.531 when 59 whole-breast
radiomic features were concatenated with 256-dimensional
MIL embeddings).
Rather than constituting a failure, this finding validates
the GatedAttentionMIL mechanism: the attention network
learns to weight axial slices by their discriminative
content, implicitly encoding the spatial information that
hand-crafted global statistics attempt to describe.
Whole-breast GLCM and shape features without a tumour mask
capture parenchymal composition rather than
tumour-specific enhancement kinetics, providing no
complementary signal beyond what the backbone already learns.

\textbf{Federated learning findings.}
The federated experiment was conducted under the extreme
non-IID partition generated by Dirichlet $\alpha=0.5$ with
seed=42, in which Hospital~1 (80 patients, Non-Luminal
dominant at 0.40:1) faces Hospitals~2 and~3 (509 patients
combined, Luminal dominant at 5.3:1 and 4.8:1 respectively).
This configuration constitutes a genuine non-IID test:
the dominant class differs across clients, not merely the
imbalance ratio.

FedAvg collapsed to trivial majority-class prediction under
this partition (F1 = 0.429, AUC = 0.565), confirming that
weighted gradient aggregation fails when dominant-class
asymmetry allows the two larger hospitals to overwhelm
the smaller minority-class hospital's gradient signal.

SCAFFOLD demonstrated an important decoupling.
Its control-variate correction successfully prevented
client drift at the representation learning level ---
reaching AUC = 0.601 by round~2, whereas FedAvg required
19 rounds to reach a comparable AUC --- but overcorrected
the decision boundary to Non-Luminal prediction (F1 = 0.200)
under the extreme imbalance gradient asymmetry.
The theoretical convergence guarantees of SCAFFOLD assume
bounded gradient dissimilarity \cite{karimireddy2020scaffold};
the extreme asymmetry of this partition places the
configuration at the boundary of that assumption, revealing
a limitation of purely drift-correcting aggregation
strategies when class imbalance compounds non-IID
heterogeneity.

FedSCRT resolved this limitation by decoupling backbone
learning from head calibration.
After SCAFFOLD backbone convergence, a federated
class-balanced classifier retraining stage aggregated local
cRT heads from all three hospitals, yielding macro F1 = 0.629
and AUC-ROC = 0.687.
The AUC exceeds the centralised baseline (0.674) by 0.013
points, consistent with the hypothesis that federated
exposure to three structurally different patient
distributions acts as an implicit regulariser.
The residual macro F1 gap of 0.038 (3.8 percentage points)
represents the cost of the extreme non-IID partition under
the 20-round communication budget.
FedSCRT closes 84.1\% of the gap that FedAvg leaves.

\section*{Limitations}

Four limitations of this work are acknowledged.

\textbf{No tumour segmentation masks.}
The dataset provides no region-of-interest annotations.
All analysis operates on the full breast volume, which
includes parenchymal tissue that carries no subtype
signal.
The performance gap relative to mask-guided published
methods is attributable in part to this constraint,
not solely to the federated setting.

\textbf{Single-site data simulation.}
The three hospital clients are simulated by partitioning
a single-institution dataset.
Real federated deployments would involve scanner differences,
acquisition protocol variations, and demographic
heterogeneity not captured by Dirichlet sampling alone.

\textbf{Communication budget.}
Twenty communication rounds may be insufficient for full
convergence under the extreme $\alpha=0.5$ partition.
The residual FedSCRT gap of 0.038 may narrow with a larger
round budget, which is left for future investigation.

\textbf{Extreme non-IID challenges for SCAFFOLD.}
The $\alpha=0.5$ partition with a reversed-dominant-class
hospital places the experiment at the boundary of SCAFFOLD's
convergence assumptions.
The joint problem of non-IID data and severe within-client
class imbalance lacks a tailored algorithm with theoretical
guarantees, representing an open research gap identified
by this work.

\section*{Perspectives and Future Work}

\textbf{Tumour segmentation masks.}
The primary enabling step for improving both the
centralised and federated results is access to tumour
segmentation masks.
The TCIA Duke-Breast-Cancer-MRI dataset provides masks
for a subset of the cohort; integrating these would
enable tumour-specific radiomics and attention
supervision, expected to improve both macro F1 and the
radiomics fusion result.

\textbf{FedSCRT on a federated SCAFFOLD backbone.}
The FedSCRT result reported in this thesis uses a
centralised SCAFFOLD backbone (all clients share the
full training set for backbone convergence) with federated
cRT heads.
Running FedSCRT on a truly federated SCAFFOLD backbone ---
where each client sees only its local partition during
all stages --- would provide a stricter privacy evaluation
and is the natural next experiment.

\textbf{Convergence-aware aggregation for joint non-IID
and class imbalance.}
SCAFFOLD's failure mode under extreme non-IID and
severe imbalance motivates the design of an algorithm
that corrects client drift while accounting for
class-conditional gradient variance.
Class-conditional control variates --- one per class per
client rather than one per client --- represent one
candidate approach that would directly address the
compound heterogeneity observed in this study.

\textbf{Differential privacy.}
The current implementation transmits model parameters
without formal privacy guarantees beyond data locality.
Integrating DP-SGD \cite{abadi2016deep} with calibrated
noise to the FedSCRT pipeline would enable a rigorous
privacy-utility trade-off analysis, which is required
for clinical deployment under Law~18-07.

\textbf{Multi-sequence and multi-modality extension.}
The current pipeline uses T1-subtraction DCE-MRI at a
single time point.
Including T2-weighted and diffusion-weighted sequences
--- which encode complementary tissue properties ---
would require a multi-branch MIL architecture but is
expected to substantially improve subtype discrimination,
particularly for the Triple-Negative class whose
T1-enhancement pattern overlaps with other subtypes.
